%==============================================================
\section{第八层：数值方法}
\label{sec:de-numerics}
%==============================================================

\subsection{有限差分法}

\subsubsection{基本格式}

\textbf{一阶导数：}
\[
u'(x_i) \approx \frac{u_{i+1} - u_{i-1}}{2h} \quad \text{（中心差分，$O(h^2)$）}
\]

\textbf{二阶导数：}
\[
u''(x_i) \approx \frac{u_{i+1} - 2u_i + u_{i-1}}{h^2} \quad \text{（$O(h^2)$）}
\]

\textbf{二维 Laplace（五点格式）：}
\[
\nabla^2 u_{i,j} \approx \frac{u_{i+1,j} + u_{i-1,j} + u_{i,j+1} + u_{i,j-1} - 4u_{i,j}}{h^2}
\]

\subsubsection{稳定性}

\textbf{热方程显式格式：} $u_j^{n+1} = u_j^n + r(u_{j+1}^n - 2u_j^n + u_{j-1}^n)$，$r = k\Delta t/h^2$

\textbf{稳定性条件：} $r \leq 1/2$

\textbf{波动方程：} $\frac{c\Delta t}{h} \leq 1$（CFL 条件）

\subsection{有限元法}

\textbf{核心步骤：}
\begin{enumerate}
  \item 将域 $\Omega$ 剖分为单元
  \item 在每个单元上用多项式基函数近似
  \item 代入弱形式，组装全局刚度矩阵
  \item 施加边界条件
  \item 求解线性系统 $K\mathbf{u} = \mathbf{f}$
\end{enumerate}

\textbf{刚度矩阵元素：}
\[
K_{ij} = \int_\Omega \nabla\phi_i \cdot \nabla\phi_j\,d\Omega
\]

\subsection{谱方法}

\textbf{思想：} 用全局正交基展开 $u(x) = \sum \hat{u}_n \phi_n(x)$

\textbf{Galerkin 方法：} 将残差投影到基函数上为零。

\textbf{优势：} 对光滑解指数收敛。

\textbf{劣势：} 对不连续解有 Gibbs 现象。

\textbf{速查版本见~\ref{subsec:pde-numerics}~节。}

\subsection{半拉格朗日方法}

对 $u_t + \mathbf{v}\cdot\nabla u = 0$：
\begin{enumerate}
  \item 从网格点 $\mathbf{x}$ 出发，沿速度场回溯：$\mathbf{x}^* = \mathbf{x} - \Delta t\,\mathbf{v}(\mathbf{x})$
  \item 在 $\mathbf{x}^*$ 处插值：$u^{n+1}(\mathbf{x}) = u^n(\mathbf{x}^*)$
\end{enumerate}

\textbf{优势：} 无条件稳定。\textbf{劣势：} 数值耗散。

%==============================================================
\section{第九层：与物理/CG 的交叉应用}
%==============================================================

\subsection{微分方程在渲染中的出现}

\begin{center}
\begin{tabular}{L{4cm}L{5cm}}
\toprule
\textbf{方程} & \textbf{渲染中的角色} \\
\midrule
Helmholtz 方程 & 波动光学、衍射 \\
辐射传输方程（RTE） & 参与介质中的光传输 \\
渲染方程（积分方程） & 全局光照 \\
扩散方程 & 次表面散射近似 \\
Poisson 方程 & 压力投影、法线重建 \\
\bottomrule
\end{tabular}
\end{center}

\subsection{微分方程在物理模拟中的出现}

\begin{center}
\begin{tabular}{L{4cm}L{5cm}}
\toprule
\textbf{方程} & \textbf{模拟对象} \\
\midrule
Navier--Stokes & 流体（水、烟、火） \\
弹性方程 & 固体变形 \\
波动方程 & 布料、绳索、声波 \\
热方程 & 温度场、扩散 \\
Hamilton 方程 & 刚体、粒子系统 \\
反应--扩散方程 & 图案生成（Turing） \\
\bottomrule
\end{tabular}
\end{center}

\subsection{微分方程在几何处理中的出现}

\begin{center}
\begin{tabular}{L{4cm}L{5cm}}
\toprule
\textbf{方程} & \textbf{几何处理任务} \\
\midrule
$\Delta u = 0$（Laplace） & 调和映射、网格平滑 \\
$\Delta^2 u = 0$（双调和） & 薄板弯曲、曲面重建 \\
平均曲率流 & 网格平滑、曲面演化 \\
测地线方程 & 最短路径、形状匹配 \\
$-\Delta u = \lambda u$ & 形状分析、谱聚类 \\
\bottomrule
\end{tabular}
\end{center}

%==============================================================
\section{总结：微分方程的统一视角}
%==============================================================

\subsection{五条核心线索}

\begin{enumerate}
  \item \textbf{对称性 → 降阶/简化}（李群方法、分离变量、行波解、自相似解）
  \item \textbf{线性算子的谱 → 正交基 → 展开}（Fourier、球谐、Bessel、Legendre）
  \item \textbf{Green 函数 → 点源响应 → 叠加}（线性问题的通用解法）
  \item \textbf{变分原理 → 弱形式 → 数值离散}（有限元的理论基础）
  \item \textbf{特征值问题 → 振动模式 → 形状/频率分析}（贯穿所有领域）
\end{enumerate}

\subsection{知识网络}

\[
\boxed{
\begin{aligned}
&\text{对称性（李群）} \xrightarrow{\text{降阶/降维}} \text{简化方程} \\
&\text{分离变量} \xrightarrow{\text{特征值问题}} \text{特殊函数（正交基）} \\
&\text{Green 函数} \xrightarrow{\text{点源叠加}} \text{卷积/积分表示} \\
&\text{变分原理} \xrightarrow{\text{弱形式}} \text{有限元/数值解} \\
&\text{Taylor/幂级数} \xrightarrow{\text{无穷维矩阵}} \text{算子对角化/积分变换}
\end{aligned}
}
\]

\subsection{从局部到全局的三条路径}

\[
\boxed{
\begin{aligned}
&\text{Taylor 展开（局部）} \xrightarrow{\text{解析延拓}} \text{复平面上的全局信息} \\
&\text{Laplace 变换（半轴）} \xrightarrow{\sigma = 0} \text{Fourier 变换（全轴）} \\
&\text{幂级数（离散指标）} \xrightarrow{\text{生成函数}} \text{复变函数（连续变量）}
\end{aligned}
}
\]

\subsection{微分方程 $\leftrightarrow$ 代数 $\leftrightarrow$ 几何的三角关系}

\[
\boxed{
\begin{aligned}
&\text{微分方程} \xleftrightarrow{\text{特征值/谱}} \text{线性代数（矩阵）} \\
&\text{微分方程} \xleftrightarrow{\text{对称性/不变量}} \text{几何（流形/群）} \\
&\text{线性代数} \xleftrightarrow{\text{表示论}} \text{几何（群作用）}
\end{aligned}
}
\]

\subsection{最终统一}

所有方法归结为一个核心问题：

\[
\boxed{Lu = f}
\]

其中 $L$ 是（微分/积分/差分）算子。求解策略：

\begin{enumerate}
  \item \textbf{找 $L$ 的特征基} → 对角化 → 逐分量求解（分离变量、谱方法）
  \item \textbf{找 $L^{-1}$（Green 函数）} → 卷积 → 叠加（Green 函数法、边界元法）
  \item \textbf{利用对称性} → 降维 → 简化（李群方法、行波解）
  \item \textbf{变分/能量最小化} → 弱形式 → 离散（有限元）
  \item \textbf{迭代/逼近} → 逐步收敛（Picard 迭代、Neumann 级数、路径追踪）
\end{enumerate}

这五条路径覆盖了从考研解题到前沿研究的所有微分方程方法。